% SPDX-License-Identifier: Apache-2.0
% covers: Timer
\datasheet{Timer}{Vreteno's core-local interruptor: a count, a compare, a software interrupt and two lines, behind an AXI peripheral end.}

\dsfacts{\code{cpu/vreteno/src/timer.rs}}{\code{vreteno32::timer::Timer}}%
{\code{//docs:vreteno}, \code{//docs:axi}}

\subsection*{Function}

\code{Timer} is the core-local interruptor of the Vreteno machine, the
CLINT. It holds a 64-bit count, \code{mtime}, a 64-bit compare,
\code{mtimecmp}, and one bit, \code{msip}, and it drives two lines. It
raises \code{tirq} while the count has reached the compare, which the
core reads as bit 7 of \code{mip}, and it raises \code{sirq} while
\code{msip} is set, which the core reads as bit 3. A program raises its
own software interrupt by writing \code{msip} and clears it by writing
it back to zero; nothing else touches that bit.

The three registers are at the offsets every RISC-V platform puts them
at, so that a stock port of an operating system finds them where it
looks. That is what makes the window 64\,KiB: \code{mtime} sits at
\code{0xbff8}, far from the compare, and the two cannot share a page.

It is the peripheral client of an \code{AxiPer}, and its link is one
port, \code{bus}, a \code{PerPort} of the four channels: requests
come on \code{bus\_req} and write beats on \code{bus\_w}, and answers
go out on \code{bus\_ans} and \code{bus\_r}. The input \code{rst}
clears the count and
nothing else.

\subsection*{Parameters}

\begin{center}\footnotesize
\begin{tabular}{@{}lp{0.7\textwidth}@{}}
\toprule
Parameter & Meaning \\
\midrule
\code{I} & The width of the AXI identifier on \code{bus\_req},
\code{bus\_ans} and \code{bus\_r}, in bits. The tables use 4, as the board
does since its bus has two hosts behind an arbiter (issue 241). \\
\bottomrule
\end{tabular}
\end{center}

\subsection*{Register map}

Offsets from the controller's base, which is
\code{isa::CLINT\_BASE}, \code{0x0200\_0000}, under the mask
\code{isa::CLINT\_MASK}, \code{0xffff\_0000}: a window of 64\,KiB.

\dsregs{Timer}

The map is declared as \code{clint}, with fourteen select bits: an
address's bits 15 to 2 select the word. A word is addressed only at its
first byte. An address with either low bit set selects no word, so it
reads zero and takes no write, as the reference model has it, and a
store of a byte or a half to a register's second lane is dropped. A
word the map does not name reads zero, and a write to it is dropped
(issue 681). The other 31 bits of a write to \code{msip} are dropped
and read back as zero.

The controller selects the word by the whole 16-bit offset, bits 15 to
0, and looks at no bit above them: the router sends it only the bursts
in its window, so nothing higher has to be decoded. The offsets are far
apart, which is why the decode is on the offset rather than on two bits
of it. Every offset in the window that is not one of the four above
reads and writes as \code{msip} does; there are five words in 64\,KiB
and nothing answers for the rest.

\dstables{Timer}

\subsection*{Behaviour}

\begin{itemize}
\item The count adds one every cycle. With \code{rst} high it goes to
zero.
\item A read is taken when \code{bus\_r} has room and no write is waiting
for its beat. It is answered on \code{bus\_r} in the cycle it is taken,
with the word addressed, \code{OKAY} and \code{last} set.
\item A write is taken when no write is waiting. The map's write
enables for the word it names, one bit a register, go into
\code{pwe}, and its identifier into \code{pid}.
\item The held write's beat is taken when it has arrived and
\code{bus\_ans} has room. The lanes the strobe covers replace those bytes of
the word, the other bytes stay, and \code{bus\_ans} sends \code{OKAY}.
\item Every cycle \code{pending} takes whether \code{mtime} is at or
above \code{mtimecmp}, compared as unsigned numbers. \code{tirq} is
\code{pending}, a register's output, so the core sees in one cycle
what the timer decided in the cycle before.
\item \code{mtimecmp} starts at zero and \code{rst} does not clear it,
so the line is high until software writes a compare above the count.
\item \code{sirq} is \code{msip} itself, with no register between, so
the core sees a program's own software interrupt in the cycle the write
lands rather than the cycle after. Nothing but a write to \code{0x0000}
changes that bit: the controller neither sets it nor clears it.
\item The count is driven every cycle, and a write to one of its halves
is driven after that, so a write in the same cycle as a tick keeps the
written word and the tick is lost.
\end{itemize}

\subsection*{Verification}

\begin{itemize}
\item \code{//cpu/vreteno:lockstep\_test} hands \code{pending} and
\code{msip} to the model for each instruction and checks
\code{mtimecmp} against the model at the halt, after draining the bus,
since a store is posted and the compare lands some cycles after the
core issued it. \code{demo\_program} sets the compare to 150 and counts
the timer's interrupt. \code{random\_programs} reads the count and
writes the compare 64 cycles beyond it, in sixty-four seeded programs.
\item \code{a\_program\_that\_interrupts\_itself} writes
\code{msip}, takes the trap, clears the bit in the handler and counts;
it is the software interrupt's own test.
\item \code{//cpu/vreteno:pair\_test} and \code{//cpu/vreteno:board\_test}
run the controller on a bus beside a core, though neither is about it.
\item The build simulates \code{Timer<2>}'s netlist against the trace
of \code{//cpu/vreteno:vreteno} under nvc and Verilator:
\code{//docs:vreteno\_sim\_timer\_timer\_tb\_test} and
\code{//docs:vreteno\_vsim\_timer\_test}.
\end{itemize}

\subsection*{Used by}

\code{Board}, as \code{timer} behind the tracker \code{ptimer}, with
\code{tirq} and \code{sirq} wired to the core's two inputs of those
names, in the 64\,KiB from \code{0x0200\_0000}. The demonstration
\code{//cpu/vreteno:vreteno}, the lockstep test and \code{run.rs}.

\subsection*{Limits}

The counter increments once per clock cycle. The controller supports
single-beat bursts only: transaction fields such as \code{len} and \code{size}
are uninspected, every read response asserts \code{last}, and multi-beat
bursts are not supported. Address decoding ignores bits above bit 15 because
the interconnect router filters traffic outside the peripheral's 64\,KiB window.
One write request is held at a time, during which incoming reads stall.
Writes update 32 bits of a 64-bit register per access, requiring two store
instructions to update a full compare threshold. The counter and compare registers
reside beyond the 12-bit immediate displacement range of a single base register,
requiring separate base pointers in software routines (\code{program.rs}).
All transactions return \code{OKAY} without generating decode errors.
